% Ad Atticum 16.15 (ad-atticum-16-15) --- English translation
% Translated by Claude (Anthropic) from the Latin (Perseus).
% Style guide: STYLE.md.

\ciceroLetterOpener{CICERO ATTICO salutem}

\ciceroSection{1} Do not think that it is from laziness that I do not write in my own hand --- though, by Hercules, it is from laziness. For I have nothing else to say. And yet in your letters too I seem to recognize Alexis. But to the matter. If Dolabella had not treated me in the most outrageous fashion, I might perhaps have been in doubt whether to be the easier-going one or to press my claim with the full rigour of the law. As things stand, I am even glad to have been handed a case in which he himself, and everyone else besides, shall see I have washed my hands of him; and I shall make a parade of it ---  \dag{}and of course, do it for my own sake and for the sake of the state\dag{}, because that is why I hate him: he had begun to defend her with my encouragement, and then, bought off with money, not only deserted her but, so far as it lay in him, actually overthrew her. As to your question, how I want the matter handled when the day comes: in the first place I should like the situation to be such that there is no impropriety in my being at Rome --- and I will do as you advise on that, as on the rest. On the substance, though, I want the case pressed thoroughly and with severity. Even so, summoning the sureties seems to carry a certain \textit{[Greek: dysōpia]} --- a hint of embarrassment; still, I wish you would consider how this would stand. For instead of summoning the sureties we can put in our agent --- they will not join issue with him. I am not unaware that, this done, the sureties are released. Yet I think it disgraceful in his case that, on a debt for which security has been given, his agents should not pay; and I think it consistent with my own dignity to pursue my right without inflicting the utmost humiliation on him. I should like you to write me back what you think on this; nor do I doubt you will handle the whole business more gently than I should.

\ciceroSection{3} I return to public affairs. Much, by Hercules, have you often written to me sensibly in the political \textit{[Greek: politikos]} line, but in this letter nothing more sensibly: for although \dag{}for the moment\dag{} for the present that boy is doing a pretty job of bouncing Antony back, still we must wait for the outcome. But what a speech he gave! For it has been sent to me. He swears ``as I hope to attain the honours of my father,'' and at the same time stretches out his right hand toward the statue. \textit{[Greek: mēde sōtheiēn hupo ge toioutou]} --- may I never be saved by the likes of him! But, as you write, I see that the most decisive moment will be the tribunate of our Casca; on which very point I said to Oppius, when he was urging me to take up the young man and his whole cause and the body of the veterans, that I could in no way do so unless I had made certain that he would not only be no enemy to the tyrant-slayers but in fact their friend. When Oppius said this would be so, I replied, ``Then why are we in a hurry? They have no need of my services before the Kalends of January, and we shall see his will toward us through Casca before the Ides of December.'' He agreed with me strongly. So much for that. As for what remains, you will have couriers every day, and, as I suppose, will every day have something to write to me about. I am sending you a copy of a letter from Lepta; from it it seems to me that that \textit{[Greek: Stratullax]} --- our strutting soldier --- has been knocked off his pedestal. But you will judge for yourself when you have read it.

\ciceroSection{4} The letter was already sealed when I received letters from you and from Sestius. Nothing more delightful than Sestius's letter, nothing more affectionate. Yours was brief; your previous letter had been generous to overflowing. You advise me, both sensibly and as a friend, to stay in these parts above all, until we have heard how these matters now stirred up will turn out.

\ciceroSection{5} But for me, my dear Atticus, the state of the republic is really not at this moment what moves me --- not that anything is dearer to me, or ought to be, but in cases beyond hope even Hippocrates forbids the application of medicine. So with that. What moves me is my private affairs. Affairs, did I say? My reputation, rather. For with such heavy obligations still outstanding to me, I have not even what is needed to discharge my debt to Terentia. Terentia, I say? You know we long ago agreed to pay off Montanus's claim of 25{,}000 sesterces. Most modestly Cicero had asked this --- on his honour. Most generously, as you yourself had approved, I had promised it, and had told Eros to keep the sum set aside. \dag{}But no\dag{}: Aurelius has been forced to borrow at ruinous interest. As for Terentia's claim, Tiro wrote to me that you said the money would come from Dolabella. I think he must have understood you badly --- if anyone ever understood anything badly --- or rather understood nothing at all. For you wrote to me what Cocceius's answer was, and Eros wrote in almost the same words. So I must come into the fire itself. It is more shameful to fall privately than publicly. And so on the other matters of which you wrote me so sweetly, with my mind so disturbed, I could not write back as I used to. \dag{}Help me in this anxiety where I am\dag{} to find a way out; I have ideas of how, but I can settle nothing for certain before I see you. And why should I not be safe where I am, when Marcellus is? But that is not the point, nor my chief concern; what my concern is, you see. So --- I am coming.
